%  08_conclusion.tex
\section{Conclusion}
\label{sec:conclusion}

We presented an open-set individual bat face recognition system using a strict
identity-disjoint evaluation protocol.
Our comparative experiments across three model families, two species, and three
background conditions lead to three conclusions:

\begin{enumerate}[leftmargin=1.6em, itemsep=3pt]
  \item \textbf{Identity-disjoint evaluation is non-negotiable.}
        The same architecture can appear to achieve $F_1 \approx 0.99$ under random-pair
        splitting and $F_1 \approx 0.74$ under identity-disjoint evaluation on the same
        images.  Any reported result on animal re-identification should explicitly state
        whether splits are identity-disjoint.

  \item \textbf{Pair-based models are the appropriate choice at the current dataset scale.}
        At $\le\!12$ training individuals, the Siamese pairwise objective consistently
        outperforms angular-margin embedding losses by $0.27$ ROC-AUC points.
        Four independent signals confirm that this is a dataset-scale limitation rather
        than a model-family limitation: embedding models reach near-perfect training
        accuracy while test ROC-AUC stays near chance; Top-1 identification is at random
        for both models; optimisation is unstable throughout training; and crucially, the
        same gap (${\approx}0.27$ ROC-AUC) appears independently for rousettus (6 train
        identities) and mauritius (5 train identities), ruling out any species-specific
        explanation and isolating identity count as the governing variable.

  \item \textbf{Green-screen or controlled backgrounds improve performance meaningfully.}
        A ${\sim}\!12$-point ROC-AUC improvement from random to controlled backgrounds
        (mauritius, Siamese) suggests that background standardisation during data
        collection is a high-leverage intervention.
\end{enumerate}

\paragraph{Reproducibility.}
All experiments are fully reproducible from the publicly released configuration files
and MLflow run IDs cited throughout this paper.  The single-command pipeline
\texttt{uv run bat-cli train --experiment <name>} regenerates the complete training,
evaluation, and reporting artefacts from scratch.
